\documentclass{standalone}
\usepackage{tikz}
\usepackage{commath}
\usepackage{newtxmath}
\tikzset{no slope/.code={\pgfslopedattimefalse}}
\renewcommand{\vec}[1]{\mathbf{#1}}
\graphicspath{{../figure_primatives/}}


\usetikzlibrary{arrows,shapes,positioning,math,calc}
\usetikzlibrary{patterns,decorations.markings,decorations.pathmorphing}
\usetikzlibrary{calc}

\begin{document}
	\begin{tikzpicture}[scale=1,every node/.style={scale=1}]
    \tikzmath{
            \panelSpacing = 10;
            \panelSpacingTwo = \panelSpacing;
            \panelOne = 5;
            \panelTwo = \panelOne + \panelSpacing;
            \panelThree = \panelTwo + \panelSpacing;
            \panelFour = 10;
            \panelFive = \panelFour+\panelSpacing;
            \verticalSpace = 10;
            \panelYTwo = 3.5;
            \colorX = 22;
            \colorY = \panelYTwo;
            \colorYTwo = \colorY+\verticalSpace;
            \panelY = \panelYTwo + \verticalSpace;
            \labelx = -3.8;
            \labely = -4;
            \boxX = 23;
            \boxY = 7;
            \heightColor=7;
        }
    \tikzstyle{boundaryStyle} = [line width=2pt,black,line cap = round];
    \tikzstyle{particleStyle} = [black!30!white];
    \tikzstyle{vectorStyle} = [line width=2pt,black,-latex,line cap= round];
    \tikzstyle{dashedStyle} = [line width=2pt,dashed];
	% \path [use as bounding box] (0,0) rectangle (\boxX cm, \boxY cm);

    \begin{scope}
        \node at (0,0) {\includegraphics{laplace_inverse_error.png}};
    \end{scope}


	\end{tikzpicture}
\end{document}